% 作者题证的编辑转录副本，不是独立下载的上游原件。
% 来源：https://github.com/BenoitDionne/PDE/blob/PDE/sobolev.tex

\paragraph{记号与假设（据作者章内约定整理）。}
$H^{k,2}=H^k$。$(k,2)$-extension property指存在有界线性延拓
$E:H^k(\Omega)\to H^k(\mathbb R^n)$，其限制回$\Omega$为恒等。
Fourier范数为$\|f\|_{k,\rho}^2=\int(1+|\xi|^2)^k|\widehat f(\xi)|^2\,d\xi$。
\begin{theorem}[Rellich's Theorem; sob\_rell\_th]
Let $0\leq k_1<k_2$ be integers. If $\Omega\subset\mathbb R^n$ is
bounded and has the $(k_2,2)$-extension property, the inclusion
$H^{k_2}(\Omega)\to H^{k_1}(\Omega)$ is compact.
\end{theorem}
\begin{proof}
Let $(f_i)$ be bounded in $H^{k_2}(\Omega)$.
Choose $\psi\in C_c^\infty(\mathbb R^n)$ equal to $1$ on $\Omega$ and
$0\leq\psi\leq1$. The operator $F=\psi E$ is bounded, satisfies
$F(f)|_\Omega=f$, and has $\supp F(f)\subset K:=\supp\psi$.
Put $g_i=F(f_i)$. Then
\[
\supp g_i\subset K,\qquad g_i|_\Omega=f_i,\qquad
\|g_i\|_{k_2,\rho}\leq C\quad\hbox{for all }i.
\]
\textbf{(i)} Choose $\phi\in C_c^\infty(\mathbb R^n)$ equal to $1$ on $K$.
Since $g_i=\phi g_i$, their Fourier transforms satisfy
$\widehat g_i=\widehat g_i*\widehat\phi$, with the normalization used for
this convolution identity. The convolution is smooth since
$\widehat\phi$ is Schwartz and $\widehat g_i\in L^2$.
The elementary weight inequality and Cauchy--Schwarz give
\begin{align*}
|\widehat g_i(y)|
&\leq\int|\widehat\phi(y-x)|\,|\widehat g_i(x)|\,dx\\
&\leq\frac{2^{k_2/2}}{(1+|y|^2)^{k_2/2}}
\int(1+|y-x|^2)^{k_2/2}|\widehat\phi(y-x)|
     (1+|x|^2)^{k_2/2}|\widehat g_i(x)|\,dx\\
&\leq2^{k_2/2}\|\phi\|_{k_2,\rho}\|g_i\|_{k_2,\rho}
\leq2^{k_2/2}C\|\phi\|_{k_2,\rho}.
\end{align*}
Differentiating the convolution, the same calculation, with
$\partial_j\widehat\phi$ in place of $\widehat\phi$, bounds
$|\partial_j\widehat g_i(y)|$ uniformly in $i,y$.
All its weighted $L^2$ norms are finite because $\widehat\phi$ is Schwartz.
Consequently the $\widehat g_i$ are uniformly bounded and equicontinuous,
the latter by the mean value theorem.

Arzel\`a--Ascoli gives a subsequence uniformly convergent on
$\overline{B_1(0)}$. Extract a further subsequence uniformly convergent
on $\overline{B_2(0)}$, and continue. The diagonal subsequence, denoted
$\widehat g_i^{[i]}$, converges uniformly on every $\overline{B_N(0)}$.

\textbf{(ii)} It suffices to prove that $(g_i^{[i]})$ converges in
$H^{k_1}(\mathbb R^n)$, because its restriction is $(f_i^{[i]})$.
Split the norm at radius $N$:
\begin{align*}
\|g_i^{[i]}-g_j^{[j]}\|_{k_1,\rho}^2
={}&\int_{|y|\leq N}(1+|y|^2)^{k_1}
 |\widehat g_i^{[i]}(y)-\widehat g_j^{[j]}(y)|^2\,dy\\
&+\int_{|y|>N}(1+|y|^2)^{k_1-k_2}(1+|y|^2)^{k_2}
 |\widehat g_i^{[i]}(y)-\widehat g_j^{[j]}(y)|^2\,dy\\
\leq{}&\sup_{|y|\leq N}|\widehat g_i^{[i]}(y)-\widehat g_j^{[j]}(y)|^2
 \int_{|y|\leq N}(1+|y|^2)^{k_1}\,dy\\
&+(1+N^2)^{k_1-k_2}\|g_i^{[i]}-g_j^{[j]}\|_{k_2,\rho}^2\\
\leq{}&\sup_{|y|\leq N}|\widehat g_i^{[i]}(y)-\widehat g_j^{[j]}(y)|^2
 \int_{|y|\leq N}(1+|y|^2)^{k_1}\,dy
 +4C^2(1+N^2)^{k_1-k_2}.
\end{align*}
Given $\epsilon>0$, first choose $N$ so that the last term is less
than $\epsilon/2$; this is possible because $k_1-k_2<0$.
For this fixed $N$, local uniform convergence supplies $J$ such that
the first term is less than $\epsilon/2$ for $i,j>J$.
Thus the diagonal sequence is Cauchy in the Banach space
$H^{k_1}(\mathbb R^n)$ and converges there. Restriction proves compactness.
\end{proof}
